\documentclass[12pt]{article}
\usepackage[utf8]{inputenc}
\usepackage{amsmath, amssymb, amsthm}
\usepackage{geometry}
\geometry{a4paper, margin=1in}

\title{Rigorous Analysis of the Erd\H{o}s-Straus Conjecture}
\author{Charles EDOU NZE\thanks{Charles EDOU NZE, chercheur indépendant}}
\date{\today}

\begin{document}
\maketitle

\begin{abstract}
This paper presents a formal analysis of the Erd\H{o}s-Straus conjecture, breaking it down into fundamental axiomatic definitions, contextual literature research, and zero-ellipse intermediate lemmas. It also provides a strict architecture for autoformalization in proof assistants.
\end{abstract}

\section{Axiomatic Definitions and Formulation}
Let $\mathbb{N} = \{1, 2, 3, \ldots\}$ be the set of natural numbers.
The Erd\H{o}s-Straus conjecture states that for every integer $n \ge 2$, there exist positive integers $x, y, z \in \mathbb{N}$ such that:
\begin{equation}
\frac{4}{n} = \frac{1}{x} + \frac{1}{y} + \frac{1}{z}
\end{equation}

To formalize this, we define the solution set for a given $n \ge 2$:
\begin{equation}
S(n) = \left\{ (x, y, z) \in \mathbb{N}^3 \;\middle|\; \frac{4}{n} = \frac{1}{x} + \frac{1}{y} + \frac{1}{z} \right\}
\end{equation}
The conjecture is equivalent to the statement: $\forall n \ge 2, S(n) \neq \emptyset$.

\section{Contextual Literature Research}
The Erd\H{o}s-Straus conjecture generalizes the Egyptian fraction representation problem. It is known that for any positive rational number $p/q$, there exist representations as sums of distinct unit fractions. The specific case of $4/n$ constraints the number of terms to three.

Historically, Mordell (1967) proved that the conjecture holds except possibly for prime numbers $n$ congruent to $1, 11, 13, 17, 19, \text{ or } 23 \pmod{24}$. Furthermore, Webb and others have established bounds using sieve methods, showing that the number of exceptions up to $N$ is at most $O(N \exp(-c \log N / \log \log N))$, which relates closely to the distribution of primes in arithmetic progressions.

By analogy with the recent resolution of weak Goldbach's conjecture by Helfgott (2013), which employed the circle method to manage additive prime combinations, one might consider a similarly localized analytic approach. However, the non-linear Diophantine nature of this equation suggests that geometric methods on algebraic surfaces, akin to those used by Manin (1986), could be more directly applicable to bounding the set of non-solutions.

\section{Lemmas and Zero-Ellipse Proof Strategy}

We decompose the problem into modular classes. A well-known identity allows us to handle many cases.

\textbf{Lemma 1 (Identity for specific modular forms):} If $n \equiv 2 \pmod{3}$, then $S(n)$ is non-empty.

\begin{proof}
Let $n \ge 2$ and assume $n \equiv 2 \pmod{3}$. This implies there exists an integer $k \ge 0$ such that $n = 3k + 2$.
We substitute $n = 3k + 2$ into the target expression $\frac{4}{n}$ and seek an expansion.
Consider the terms $x = n$, $y = k+1$, and $z = n(k+1)$. We must verify that:
\begin{equation}
\frac{1}{n} + \frac{1}{k+1} + \frac{1}{n(k+1)} = \frac{4}{n}
\end{equation}
Starting from the left side, we find a common denominator $n(k+1)$:
\begin{equation}
\frac{1}{n} + \frac{1}{k+1} + \frac{1}{n(k+1)} = \frac{k+1}{n(k+1)} + \frac{n}{n(k+1)} + \frac{1}{n(k+1)}
\end{equation}
Summing the numerators, we obtain:
\begin{equation}
\frac{k+1 + n + 1}{n(k+1)} = \frac{n + k + 2}{n(k+1)}
\end{equation}
Substituting $n = 3k + 2$ into the numerator gives:
\begin{equation}
\frac{3k + 2 + k + 2}{n(k+1)} = \frac{4k + 4}{n(k+1)} = \frac{4(k+1)}{n(k+1)}
\end{equation}
Since $k \ge 0$, $k+1 \neq 0$, we can cancel $k+1$ from the numerator and denominator:
\begin{equation}
\frac{4(k+1)}{n(k+1)} = \frac{4}{n}
\end{equation}
Thus, $(n, k+1, n(k+1)) \in S(n)$, proving the lemma.
\end{proof}

\section{Architecture for Autoformalization}
To integrate this proof into an assistant like Lean 4, the structure must be explicitly typed:

\begin{verbatim}
import Mathlib.Data.Nat.Basic
import Mathlib.Data.Rat.Basic

-- Axiomatic Definitions
def S (n : Nat) : Set (Nat * Nat * Nat) :=
  { p | 4 * (p.1 * p.2.1 * p.2.2) =
        n * (p.2.1 * p.2.2 + p.1 * p.2.2 + p.1 * p.2.1)
        /\ p.1 > 0 /\ p.2.1 > 0 /\ p.2.2 > 0 }

-- Conjecture Type
def erdos_straus_conjecture : Prop :=
  forall n : Nat, n >= 2 -> Not (S n = empty)

-- Lemma 1 Type
lemma lemma_mod3 (n : Nat) (h1 : n >= 2) (h2 : n % 3 = 2) : Not (S n = empty) :=
  sorry
\end{verbatim}

\end{document}
